% =====================================================================
%  FICHE 11 — THEME 5 — Corrigé DIFFICILE — Piles, Nernst, Faraday
% =====================================================================
\documentclass[11pt,a4paper]{article}
\usepackage[utf8]{inputenc}
\usepackage[T1]{fontenc}
\usepackage{lmodern}
\usepackage[french]{babel}
\usepackage{amsmath,amssymb}
\usepackage[version=4]{mhchem}
\usepackage{siunitx}
\sisetup{output-decimal-marker={,},per-mode=symbol}
\usepackage{xcolor}
\usepackage{array}
\usepackage{booktabs}
\usepackage{enumitem}
\usepackage[most]{tcolorbox}
\usepackage[a4paper,margin=2.0cm,headheight=26pt,headsep=10pt]{geometry}
\usepackage{fancyhdr}
\usepackage[hidelinks]{hyperref}

\definecolor{primary}{RGB}{146,95,160}
\definecolor{primarydark}{RGB}{92,52,104}
\definecolor{corrcol}{RGB}{0,135,90}
\newcommand{\NO}{\ensuremath{\mathrm{N.O.}}}
\newcommand{\Ered}{\ensuremath{E_{\mathrm{r}}^{\circ}}}
\newcommand{\fem}{\ensuremath{\Delta E}}
\newcommand{\cc}[1]{\left[#1\right]}

\newcounter{exo}
\newtcolorbox{corr}[1][]{enhanced,breakable,boxrule=0.9pt,arc=2pt,colback=corrcol!4,
  colframe=corrcol,fonttitle=\bfseries,coltitle=white,left=3mm,right=3mm,top=2.2mm,bottom=2.2mm,
  title={Corrigé~\refstepcounter{exo}\theexo\ifx\relax#1\relax\relax\else\ \textmd{--- \itshape #1}\fi}}

\pagestyle{fancy}\fancyhf{}
\fancyhead[L]{\small\textcolor{primarydark}{\textbf{Fiche 11 · Thème 5} — Corrigé}}
\fancyhead[R]{\small\textcolor{primarydark}{\textsc{Difficile}}}
\fancyfoot[C]{\small\thepage}
\renewcommand{\headrulewidth}{0.8pt}
\renewcommand{\headrule}{\hbox to\headwidth{\color{primary}\leaders\hrule height \headrulewidth\hfill}}

\begin{document}
\begin{center}
{\LARGE\bfseries\color{primarydark}Thème 5 — Corrigé Difficile}\\[-2pt]
{\color{corrcol}\rule{\textwidth}{1pt}}
\end{center}

\begin{corr}[analyser une pile cuivre--nickel]
\begin{enumerate}[label=\textbf{\alph*)},leftmargin=2.2em,topsep=2pt,itemsep=2pt]
  \item $\Ered(\ce{Ni^{2+}}/\ce{Ni})=-0{,}26 < \Ered(\ce{Cu^{2+}}/\ce{Cu})=+0{,}34$ :
        le \textbf{nickel est l'anode} ($-$), le \textbf{cuivre est la cathode} ($+$).
  \item Anode (Ni) : $\ce{Ni -> Ni^{2+} + 2 e^-}$ (oxydation).
        Cathode (Cu) : $\ce{Cu^{2+} + 2 e^- -> Cu}$ (réduction).
  \item $\fem^{\circ}=\Ered(\text{cathode})-\Ered(\text{anode})=0{,}34-(-0{,}26)=\boxed{0{,}60\ \mathrm{V}}$.
  \item Électrons : de l'\textbf{anode (Ni) vers la cathode (Cu)} dans le fil. Pont salin : les \textbf{anions
        \ce{Cl^-}} migrent vers l'anode (compartiment \textbf{Ni}), les \textbf{cations \ce{K+}} vers la
        cathode (compartiment \textbf{Cu}).
  \item \textbf{(A)} faux (à l'électrode de Cu = cathode, c'est la \emph{réduction} \ce{Cu^{2+} -> Cu}) ;
        \textbf{(B)} faux (Ni est l'anode) ;
        \textbf{(C)} \textbf{vrai} (électrons Ni $\to$ Cu) ;
        \textbf{(D)} faux (les \ce{Cl^-} vont vers le compartiment du \emph{nickel}, pas du cuivre).
        \emph{Réponse C (annales 2020-Q18).}
\end{enumerate}
\end{corr}

\begin{corr}[f.é.m. d'une pile hors conditions standard]
\begin{enumerate}[label=\textbf{\alph*)},leftmargin=2.2em,topsep=2pt,itemsep=2pt]
  \item $-0{,}76<+0{,}80$ : anode = \textbf{zinc}, cathode = \textbf{argent}.
  \item Cathode (\ce{Ag+}/\ce{Ag}, $n=1$) :
        \[
        E_{\text{cath}} = 0{,}80 + 0{,}06\,\log(0{,}10) = 0{,}80 + 0{,}06\times(-1) = \SI{0.74}{V}.
        \]
  \item Anode (\ce{Zn^{2+}}/\ce{Zn}, $n=2$) :
        \[
        E_{\text{anode}} = -0{,}76 + \frac{0{,}06}{2}\,\log(0{,}010) = -0{,}76 + 0{,}03\times(-2) = \SI{-0.82}{V}.
        \]
  \item $\fem = E_{\text{cath}} - E_{\text{anode}} = 0{,}74 - (-0{,}82) = \boxed{1{,}56\ \mathrm{V}}$.
\end{enumerate}
Remarque : on retrouve la valeur standard ($1{,}56$\,V) car la chute de $0{,}06$ à la cathode est
\emph{exactement} compensée par celle à l'anode. \emph{(Réponse A, livre QCM~26.)}
\end{corr}

\begin{corr}[affinage électrolytique du cuivre]
$n=2$, $M(\ce{Cu})=63{,}5$.
\begin{enumerate}[label=\textbf{\alph*)},leftmargin=2.2em,topsep=2pt,itemsep=2pt]
  \item $t = 1{,}0\ \mathrm{h} = \SI{3600}{s}$ :
        \[
        m = \frac{M\,I\,t}{n\,F} = \frac{63{,}5\times 5{,}0\times 3600}{2\times 96485}
        = \frac{1\,143\,000}{192\,970} \simeq \boxed{5{,}92\ \mathrm{g}}.
        \]
  \item On isole $t$ pour $m=\SI{20}{g}$ :
        \[
        t = \frac{m\,n\,F}{M\,I} = \frac{20\times 2\times 96485}{63{,}5\times 5{,}0}
        = \frac{3\,859\,400}{317{,}5} \simeq \SI{12156}{s} \simeq \boxed{3{,}4\ \mathrm{h}}.
        \]
\end{enumerate}
La masse déposée est proportionnelle à la charge $Q=It$ : pour un dépôt donné, plus l'intensité est
forte, plus l'électrolyse est rapide.
\end{corr}

\end{document}
